\section{Round Trip}
\label{sec:cses-round-trip}

\problemstatement{Given an undirected graph of $n$ nodes and $m$ edges,
find any cycle (a route that starts and ends at the same node using
distinct edges) and print the sequence of nodes, or \texttt{IMPOSSIBLE}.}

\begin{patternbox}
\emph{``Find a cycle in an undirected graph''} is a \textbf{DFS} that
watches for a \textbf{back edge}: an edge to an already-visited node that
is not the immediate parent. Such an edge closes a cycle, which is
reconstructed from the DFS parent chain.
\end{patternbox}

\subsection*{Back edge closes a cycle}

Run DFS tracking each node's parent. If DFS at node $u$ finds a neighbour
$v$ that is already visited and $v$ is not $u$'s parent, then $v$ is an
ancestor of $u$ and the tree path from $v$ down to $u$, plus the edge
$u\to v$, forms a cycle. Reconstruct it by following parents from $u$ up to
$v$.

\begin{ideabox}[Not-the-Parent Is the Whole Condition]
In an undirected graph every edge appears from both ends, so seeing the
parent again is not a cycle. A visited neighbour that is \emph{not} the
parent is a genuine back edge -- the signal of a cycle. Recording parents
lets you print it.
\end{ideabox}

\subsection*{Algorithm}

\begin{enumerate}
  \item DFS from each unvisited node, storing parents.
  \item On finding a visited non-parent neighbour $v$ from $u$: walk
        parents $u\to\cdots\to v$, append $v$ again to close the loop, and
        print.
  \item If no back edge is ever found, print \texttt{IMPOSSIBLE}.
\end{enumerate}

\subsection*{Dry run}

\dryrun{Edges $1\!-\!2$, $2\!-\!3$, $3\!-\!1$.}

\begin{itemize}
  \item DFS $1\to2\to3$; from $3$, neighbour $1$ is visited and not
        $3$'s parent ($2$) -- back edge.
  \item Cycle: $1\to2\to3\to1$.
\end{itemize}

\subsection*{C++ implementation}

\begin{lstlisting}
#include <bits/stdc++.h>
using namespace std;

int n, m;
vector<vector<int>> adj;
vector<int> parent;
vector<bool> vis;
int cycleStart = -1, cycleEnd = -1;

bool dfs(int u, int p) {
    vis[u] = true;
    parent[u] = p;
    for (int v : adj[u]) {
        if (v == p) continue;                    // skip the edge back to parent
        if (vis[v]) { cycleStart = v; cycleEnd = u; return true; }
        if (dfs(v, u)) return true;
    }
    return false;
}

int main() {
    cin >> n >> m;
    adj.assign(n + 1, {});
    for (int i = 0; i < m; i++) {
        int a, b; cin >> a >> b;
        adj[a].push_back(b);
        adj[b].push_back(a);
    }
    parent.assign(n + 1, 0);
    vis.assign(n + 1, false);

    for (int s = 1; s <= n; s++)
        if (!vis[s] && dfs(s, -1)) break;

    if (cycleStart == -1) { cout << "IMPOSSIBLE\n"; return 0; }
    vector<int> cyc;
    for (int v = cycleEnd; v != cycleStart; v = parent[v]) cyc.push_back(v);
    cyc.push_back(cycleStart);
    cyc.push_back(cycleEnd);                      // close the loop
    reverse(cyc.begin(), cyc.end());

    cout << cyc.size() << "\n";
    for (int v : cyc) cout << v << " ";
    cout << "\n";
    return 0;
}
\end{lstlisting}

\begin{complexitybox}
\textbf{Time Complexity.} $O(n+m)$ -- one DFS. \textbf{Space Complexity.}
$O(n+m)$, plus recursion depth $O(n)$.
\end{complexitybox}

\begin{takeawaybox}
Undirected cycle detection is a DFS back edge to a non-parent ancestor;
parents reconstruct the cycle. (Disjoint Set also detects a cycle -- the
first edge joining two already-connected nodes -- but does not print it as
directly.)
\end{takeawaybox}
